% FRAGMENT A09 (drafter A: class model and main results)
% REPLACES: main.tex lines 691–693 ("The benchmark uses incumbent strengths" to "… gross acquisition profit at")
% CHANGE: states that Sections 4.2 and 4.3 are the point-mass benchmark (prior 1/2, money unit
%   l=1). Every number in Sections 4.2 and 4.3 is unchanged: with the low cost at zero, the
%   low-cost type still prepares at every price, as it did at the old low cost 1, which lay
%   below gross profit at the lowest feasible belief at every strength used there
%   (B_r(m) = 2.728490, 2.366179, 2.272133 at r = 1.2, 3, 3.6; 2.61 to 2.64 at the certified
%   strengths 1.55 to 1.65). Only the net-surplus levels inside Table 2 move (pipeline output).
% LABELS: none kept, removed, or added.  NUM PLACEHOLDERS: none.
The point-mass benchmark has prior \(1/2\), measures money in units of
the low value, \(\ell=1\), and uses incumbent strengths \(r_0=1.2\),
\(r_1=3\), and \(r_2=3.6\); \hyperref[pa-parameters]{Appendix A.8} lists
the full parameter vectors. The numbers in this section and in
\hyperref[sec-coexistence-and-the-equilibrium-correspondence]{Section 4.3}
refer to this benchmark. From \(r_0\) to \(r_1\), gross acquisition profit at
